\section{Independent Provenance Replay}

\subsection{Four propositions, not one slogan}

The story requires four propositions: table napkins were customary in the
relevant Jewish setting; disposal communicated whether a master had finished;
folding and crumpling carried the alleged opposite meanings; and John's
\emph{soudarion} and participle invoke that code. Evidence for one cannot
establish all four. A Roman \emph{mappa} proves dining cloths existed, not that
Palestinian masters folded them as messages.

\subsection{The directly relevant rabbinic dispute}

Mishnah Berakhot 8:3 records Shammai and Hillel disagreeing over where a cloth
used to wipe hands should be placed during a meal: table or cushion. Babylonian
Talmud Berakhot 52b explains the dispute through ritual impurity passing among
moisture, cloth, table, food, and cushion. There is no departing master, servant
awaiting a signal, end-of-meal fold, or promise of return.

This is a material negative because the text genuinely concerns a cloth at a
meal. It shows the danger of keyword apologetics: a ``Jewish napkin custom''
exists, but not the alleged one.

Mishnah Sanhedrin 6:1 supplies an actual cloth signal: a person waves one to
halt an execution when exculpatory evidence appears. Rabbinic literature can
state signals plainly. This signal is unrelated to dining or return.

\subsection{Greco-Roman controls}

Roman diners could bring a \emph{mappa}, use it, and carry food away in it.
Servants could hold towels during washing. Smith's antiquities article gathers
such functions without a finished/returning code. The famous ancient public
signal involving a dropped cloth begins circus races, not a domestic promise.
Cloth gestures were possible; analogy is not attestation.

\subsection{Lightfoot's bounded silence}

Lightfoot's \emph{Horae Hebraicae et Talmudicae} is useful because it pursues
rabbinic Gospel parallels verse by verse. In the four-volume OCR searched for
the source study, ``napkin'' occurs only at Lazarus, and John 20 has no verse-7
note. OCR misses and Lightfoot's selection remain possible. His silence is not
universal proof; it is evidence that an eager early-modern collector did not
transmit the custom.

\subsection{The located modern trail}

The earliest dated online presentation located is a reader's question posted
by \emph{Jerusalem Perspective} in October 2006. David Bivin requested the
missing evidence and suggested the archaic English ``napkin'' helped generate
the story. An archived email forward appears in May 2007. Later assessments and
apologetic pages repeat a 2007 circulation estimate without supplying ancient
evidence.

This is a lower bound, not an origin certificate: earlier oral or print
circulation could emerge. The publishable conclusion is ``unattested in the
enumerated corpus and routes searched,'' not ``no Jew ever did this.'' A
synonym, untranslated source, damaged OCR, or unrecorded custom could escape.

The cumulative judgment remains strong. Four propositions lack joint support;
the located trail is late; the Greek lexical control points to wrapping; and
commentators positioned to transmit rabbinic parallels do not know the story.

